% Notes of the Montreal seminar, 1970 (cote 129, typed part): archive cote 129, site page 4 (PDF page 5).
% typescript. Draft compared: none.
% AI-transcribed by Claude Opus 5.5 (claude-opus-5-5, adaptive thinking, effort high), Message Batches API, batch mtl129_full (msgbatch_01AGnDKHwy6ATCM7ehUa6263), 2026-11-13; not yet checked by eye.
\scanpage{129/4}
Par passage à la limite, les morphismes $V_n$ définissent
un morphisme additif
\[ V\colon \mathbb{W} \to \mathbb{W} \]
appelé \uline{Verschiebung} ou \uline{morphisme de décalage}.

2.2. Le \uline{morphisme de frobénius} $F\colon \mathbb{W} \to \mathbb{W}$ est défini par passage
à la limite à partir des morphismes
\[ F_n\colon \mathbb{W}_n \to \mathbb{W}_n \qquad (x_1,x_2,\dots,x_n) \mapsto (x_1^p,x_2^p,\dots,x_n^p). \]

Ce morphisme n'est en général par additif, mais par réduction
en caractéristique $p$, il définit un \uline{homomorphisme d'Anneau}
\[ F_{\mathbb{F}_p} : \mathbb{W}_{\mathbb{F}_p} \to \mathbb{W}_{\mathbb{F}_p} \qquad (\text{resp. } F_{n\,\mathbb{F}_p}) \]
où $\mathbb{F}_p$ désigne le corps premier $\mathbb{Z}/p\mathbb{Z}$ et $\mathbb{W}_{\mathbb{F}_p}$ désigne le schéma
induit sur $\mathbb{F}_p$ par changement de base. En effet les deux lois de
composition sur $\mathbb{W}$ sont données par
\[ (x_1,x_2,\dots,x_n,\dots) + (y_1,\dots,y_n,\dots) = (S_1,S_2,\dots,S_n,\dots) \]
\[ (x_1,x_2,\dots,x_n,\dots)\cdot(y_1,\dots,y_n,\dots) = (P_1,\dots,P_n,\dots) \]
où les $S_i$ et $P_i$ sont des polynômes à coefficients entiers en
$(x_1,x_2,\dots,x_i,\ y_1,y_2,\dots,y_i)$. Or pour toute $\mathbb{F}_p$-algèbre, on a
\[ P(x^p,y^p) = [P(x,y)]^p. \]

\mspage{\hand{4}}
